\subsection{Complete quadratic module and construction domains}

For the six residual indices $1055,1768,1795,3258,3413,3472$, connected
rank-two based modules admit a complete finite classification. The
positive Frobenius--Perron character and its quadratic conjugate exhaust
their two-dimensional rational representation. If a self-dual generator
has dimension polynomial $t^2-Tt+U$, its action is
\[
 \begin{pmatrix}a&b\\b&T-a\end{pmatrix},\qquad
 0\leq a\leq T,\qquad b^2=a(T-a)-U>0.
\]
Vertex interchange allows $a\geq T-a$. All other action matrices follow
from rational affine expressions in the generator dimension. Exact
nonnegative integrality, duality, and all fusion products leave respectively
$1,1,2,1,1,2$ based-module classes. Their full actions and candidate
internal-end objects are recorded. These data do not solve module
associativity and are not constructions of categories.

For index 1055 the forced $1+g$-module difference representation has rank
two. Normality and the source annihilator reduce the first generator to
zero; all 1296 remaining bounded simultaneous possibilities give exactly
four representations, each with $S_g=-I$. In a central Tannakian
condensation the tensor quotient sends $g$ to the unit, which would instead
force $S_g=I$. Thus this condensation route is impossible.

An exact search through all 128 unit-containing subsets shows that the
only proper nontrivial fusion subring of indices 1055,1768,1795 is the
pointed $C_2$ subring. Indices 3258,3413,3472 have no proper nontrivial
fusion subring. The irrational total dimensions therefore rule out every
nontrivial faithful finite-group grading: its neutral component would
have integral dimension one or two. A nontrivial equivariantization
contains $\operatorname{Rep}(G)$; the same subring audit excludes it for
the latter three rings and leaves only $G=C_2$ for the former three.

For 1768 and 1795 the sign invertible acts freely on all eight simples.
The orbit/projective-stabilizer description of equivariantization simples
then forces a hypothetical $C_2$-de-equivariantization to have rank four:
a free base orbit would give a sign-fixed simple, while a fixed base simple
gives two equivariant lifts. Its dimension would be respectively
$222+36\sqrt{37}$ or $270+120\sqrt5$, both greater than $45$.
The unrestricted rank-four companion theorem bounds that dimension by
$24+12\sqrt3<45$. This explicitly cited theorem dependency rules out both
routes; its source digest and the exact quotient tensors are pinned in
the evidence. Together with the $S_g$ argument above, all six residuals
are excluded from nontrivial finite-group equivariantization constructions.
No central lift is inferred from commutativity of the fusion ring alone.

There is also an unbounded arithmetic sieve for generalized-near-group
construction towers. If a base has $G$ invertibles, stabilizer order $s$,
noninvertible orbit size $\ell=G/s$, common irrational dimension $d$, and
total noninvertible multiplicity $K$, then
\[
 d^2=s+Kd,\qquad D_0=\ell(2s+Kd),\qquad
 \frac{\operatorname{N}(D_0)}{\operatorname{Tr}(D_0)}=G.
\]
Pointed products, finite-group extensions, and equivariantizations multiply
dimension by a positive integer $k$, while Morita equivalence preserves it.
Consequently $G_D=\operatorname{N}(D)/\operatorname{Tr}(D)$ must be integral,
$k\mid G_D$, $s\mid G_D/k$, and
$K^2=\operatorname{Tr}(D)/(k\ell)-4s$. These divisor and square tests
exhaust the stated family without a parameter cutoff. Division by
de-equivariantization is outside its scope.

For 1795,3258,3413, the values $G_D=10/3,45/7,45/7$ exclude the whole
stated family. For 1055 and 1768, $G_D=6$ and all matches have $s=1,K=12$;
the tuples $(k,G,\ell)$ are
\[
 (1,6,6),\quad (2,3,3),\quad (3,2,2),\quad (6,1,1).
\]
The last is excluded by the rank-two classification and the third by the
rank-four bound, its base dimension being $148+24\sqrt{37}>45$.
For 3472, $G_D=7$, $s=1,K=14$, and the matches are $(1,7,7),(7,1,1)$;
the latter is again rank-two impossible. The remaining matches are
dimension constraints only, without an asserted base category or Morita
equivalence.

\subsection{Exact character-field coverage}

For the seven indices below, the source element
$H=\sum_{j=0}^7(j+1)^2X_j$ has squarefree degree-eight characteristic
polynomial and generates the rational algebra. Exact interpolation writes
every basis element as a rational polynomial in $H$. Factoring its
polynomial over the quadratic FP field therefore proves the following
full-character coverage, independently of numerical root approximations.

\begin{center}
\begin{tabular}{r|r|r|l}
Index & FP radicand & Characters in FP field & Other factor degrees\\\hline
65 & 5 & 4 & $2,2$\\
1055 & 37 & 6 & $2$\\
1768 & 37 & 8 & none\\
1795 & 5 & 8 & none\\
3258 & 5 & 8 & none\\
3413 & 5 & 5 & $3$\\
3472 & 2 & 2 & $3,3$
\end{tabular}
\end{center}

Thus a modular-eigenvalue calculation confined to the real FP field has
complete character coverage for 1768,1795,3258. The other four need larger
splitting fields. The following exact adapters now supply all four of
them, with every source character polynomial checked to split completely:
\[
\begin{array}{c|c|l}
\text{index}&[L:\mathbb Q]&\text{defining polynomial of }u\\\hline
65&4&u^4-8u^2+36\\
1055&4&u^4-68u^2+1600\\
3413&6&u^6-21u^4+2u^3+84u^2+24u-19\\
3472&6&u^6+2u^5-9u^4-14u^3+10u^2+8u+1
\end{array}
\]
For 65 and 1055, the fields are respectively
$\mathbb Q(\sqrt5,i)$ and $\mathbb Q(\sqrt{37},\sqrt{-3})$, with
$u$ the sum of the displayed radicals. Complex conjugation fixes the
real radical and negates the imaginary one; its multiplicativity and
involution identities are checked exactly. For 3413 and 3472, take
$u=\sqrt q+\theta$, with respectively
$(q,p_\theta)=(5,t^3-3t+1)$ and $(2,t^3+t^2-2t-1)$.

For $\lambda\in L$, algebraic integrality is equivalent to the rational
matrix of multiplication by $\lambda$ having monic integral characteristic
polynomial: necessity follows from the minimal polynomial, and sufficiency
because that characteristic polynomial annihilates $\lambda$.
This tests the full ring of integers without assuming that the primitive
power basis is integral. The adapters check all eight character rows,
their products, dualities, orthogonality, the FP embedding, and direct
multiplication-operator controls. All twelve distinct restriction columns
of the known center of Fibonacci $\boxtimes\operatorname{Vec}(C_2\times C_2)$
are retained in both overfield adapters. These are complete arithmetic
applicability results; categorical obstruction claims retain the
hypotheses of the associated center theorem.

% Follow this fragment by connected_algebra_section.tex, which contains
% the general connected-algebra proof and the complete enumeration table.
